\section{Podsumowanie i wnioski}
\label{sec:podsumowanie}

Rozdział odpowiada na pytania ze wstępu w tej samej numeracji, a następnie podaje wnioski ogólne i ograniczenia przeprowadzonych badań.

\subsection{Odpowiedzi na pytania}
\label{sec:odpowiedzi}

\begin{enumerate}
    \item Reguła Oji w wariancie podprzestrzeni stabilizuje wszystkie 74 wzorce dostarczonych zbiorów, reguła Hebba 12 z nich, a wariant jednoneuronowy żadnego.
    \item Przy regule Oji sieć uczy się poprawnie wszystkich zbiorów. Przy regule Hebba żaden zbiór nie jest zapamiętany w całości, a większość wzorców jest stabilna tylko w \texttt{large-25x25} i \texttt{small-7x7}.
    \item Wzorce nie są stabilne, gdy suma ich nakładań z pozostałymi wzorcami jest duża w porównaniu z sygnałem. Dzieje się tak w zbiorach ze wspólnym tłem lub wspólną częścią, nawet przy małym obciążeniu $p/N$. W zbiorze liter niestabilne są wyłącznie piksele, w których litera różni się od większości pozostałych liter. Reguła Oji usuwa ten efekt, bo zbiega do rzutu na podprzestrzeń wzorców.
    \item Wspólna dolna połowa bitmap \mbox{25x50} zwiększa wszystkie nakładania zgodnie z~\eqref{eq:nakladanie-25x50}, więc przy regule Hebba żaden wzorzec nie jest stabilny i sieć nie odzyskuje żadnego wzorca. Przy regule Oji oba zbiory są odzyskiwane tak samo dobrze.
    \item Przy regule Oji sieć odzyskuje prawie wszystkie wzorce dużych i mało skorelowanych zbiorów do szumu od \num{0.35} do \num{0.4} i do zasłonięcia od \num{0.3} do \num{0.45}. Zbiory małe i skorelowane mają mniejsze obszary przyciągania, a najmniejszy ma \texttt{letters-14x20}. Wraz z liczbą wzorców skuteczność spada powoli dla obrazów i szybko dla liter. Przy regule Hebba skuteczność jest ograniczona odsetkiem stabilnych wzorców i tylko w \texttt{large-25x25} przekracza połowę przy umiarkowanym szumie.
    \item Przy regule Oji zbiór zróżnicowany jest odzyskiwany lepiej niż zbiór z grupami bardzo podobnych wzorców. Poziom krytyczny $f_{0,9}$ (rozdział~\ref{sec:metryki}) wynosi dla nich \num{0.4} wobec \num{0.25} dla $d = \num{0.02}$, a przy mniejszym podobieństwie ($d = \num{0.2}$) oba zbiory wypadają tak samo. Błędy w zbiorach grupowych zostają w obrębie grupy. Reguła Hebba nie odzyskuje żadnego wariantu, bo zlewa każdą grupę w jeden stan, a w zbiorze zróżnicowanym odzyskuje tylko jeden stabilny wzorzec.
    \item Losowe wejście kończy w jednym z minimów lokalnych energii. Przy regule Oji wzorce są minimami globalnymi, więc sieć zwykle odtwarza jeden z wzorców albo jego negację, każdy z podobną szansą, a pozostałe wejścia trafiają do wielu płytkich stanów pozornych. Przy regule Hebba najgłębsze są stany pozorne, często mieszaniny kilku wzorców, i to w nich kończy prawie każde losowe wejście. Wyjątkiem są silnie skorelowane litery \texttt{letters-14x20}, w których przy obu regułach sieć prawie nigdy nie trafia do wzorca, a w trybie synchronicznym przy regule Oji zwykle oscyluje.
    \item Zbiór $S_u$ wszystkich bitmap \mbox{5x5}, w których suma pierwszych 24 pikseli jest równa $-2 \xi_{25}$, ma \num{4992288} elementów. Każdy z nich jest ściśle stabilny zarówno w sieci nauczonej regułą Hebba na całym zbiorze, jak i w sieci nauczonej regułą Oji na jego bazie. Konstrukcja wynika z reguły rzutowej, która stabilizuje każdy wektor podprzestrzeni wzorców, a liczność zbioru jest bliska granicy dla hiperpłaszczyzn.
    \item Zbiór uczący złożony z kształtów T, L, H i + (reguła Hebba) i wejście w postaci szachownicy, która ma małe nakładanie z każdym kształtem. Dzięki wyzerowanej przekątnej pole ma w każdym pikselu znak przeciwny do stanu, więc w trybie synchronicznym sieć oscyluje między szachownicą a jej negacją.
    \item Na zdjęciach o \num{60000} pikseli reguła Oji stabilizuje i odzyskuje wszystkie wzorce aż do szumu \num{0.45}, a przy zasłonięciu traci je dopiero przy kilkudziesięciu procentach zakrytych wierszy. Reguła Hebba mimo bardzo małego obciążenia stabilizuje tylko zdjęcia o małej sumie nakładań z pozostałymi, a pozostałe zdjęcia kończą w stanach pozornych. Sieć tej wielkości wymaga wag w postaci czynnikowej. Zajmują one kilka megabajtów zamiast \qty{26.8}{GiB}, a uczenie i odtwarzanie trwają milisekundy.
\end{enumerate}

\subsection{Wnioski i ograniczenia}
\label{sec:wnioski}

O skuteczności reguły Hebba decydują korelacje wzorców, a nie ich liczba. Reguła ta zawodzi na zbiorach ze wspólnym tłem, na bitmapach \mbox{25x50} i na zdjęciach kotów, choć obciążenie $p/N$ wynosi dla kotów tylko $2 \cdot 10^{-4}$. Reguła Oji zbiega do rzutu na podprzestrzeń wzorców (rozdział~\ref{sec:regula-oji}), więc usuwa wpływ korelacji na stabilność i stabilizuje wszystkie wzorce dostarczonych zbiorów i zdjęć kotów. Korelacje nadal ograniczają jednak odtwarzanie. Silnie skorelowane litery mają najmniejsze obszary przyciągania, a część losowych wejść kończy w jednym z wielu płytkich stanów pozornych.

Badania mają kilka ograniczeń. Reguła Oji jest liczona wsadowo na całym zbiorze wzorców (rozdział~\ref{sec:uczenie}), a nie wzorzec po wzorcu jak reguła Hebba. Gwarancja stabilności z rozdziału~\ref{sec:regula-oji} dotyczy wzorców liniowo niezależnych, więc ich liczba nie może przekroczyć $N$. Czasy z rozdziału~\ref{sec:czas-pamiec} zmierzono na jednym komputerze, więc ich wartości bezwzględne zależą od sprzętu.
